\begin{abstract}
Weight-only post-training quantization (PTQ) costs some recognition datasets far more accuracy than others.
We ask whether the dataset affects quantization robustness, and how.
With half-quadratic quantization (HQQ), we quantize three backbones, each trained with three training strategies on \V{nDatasets} fine-grained datasets, at \V{nConditions} combinations of bit width and group size.
Each quantized model is compared with its own checkpoint through \term{accuracy retention}, the ratio of quantized to baseline accuracy.
Bit width and group size set the level of retention, and the datasets separate where retention lies between collapse and full retention.
There, five ultra-fine-grained cultivar datasets (UFG) lose more than the \V{nOther} fine-grained datasets (FG): at 3-bit weights with group size 128 they retain \V{retSc_w3g128} of their accuracy, against \V{retOther_w3g128}.
The \term{vulnerable-tail margin} of the baseline logits correlates with this difficulty and separates some datasets that baseline accuracy cannot.
Matching the margin distributions of the two groups narrows their \term{retention gap} by \V{reductionMin}\% to \V{reductionMax}\%, without closing it.
Feature similarity between the baseline and quantized models does not track the accuracy loss.
\end{abstract}
